% ============================================================
% Section 2: Related Work
% BearMamba-3: Mamba-3 + L_kin + Multi-Sensor Fusion
% Target: MSSP (IF 8.4)  Draft: 2026-06-15
% ============================================================

\section{Related Work}
\label{sec:related}

\subsection{Classical Signal Processing Baselines}
\label{subsec:rel_classical}

Classical bearing diagnosis pipelines combine signal processing with shallow
classifiers: envelope analysis demodulates the carrier to expose
fault-frequency sidebands, spectral kurtosis (SK) identifies the optimal
demodulation band, and hand-crafted statistical features (RMS, kurtosis,
crest factor) feed linear SVMs or k-nearest-neighbour
classifiers~\cite{antoni2006spectral_kurtosis,randall2011rolling}.
These methods are interpretable and require no GPU, but their feature sets
are fixed at design time and do not generalise to variable shaft speed without
explicit per-sample recomputation of fault frequencies.
We include a representative SK-SVM baseline in our CWRU evaluation
(Table~\ref{tab:cwru_baselines}) to provide an absolute accuracy anchor
and to highlight where deep models, and BearMamba-3 in particular, add value
beyond hand-crafted features.

\subsection{Deep Learning for Bearing Fault Diagnosis}
\label{subsec:rel_dl}

Convolutional neural networks (CNNs) established a dominant paradigm for
bearing fault diagnosis: 1-D convolutions extract periodic impulse features
directly from raw vibration waveforms, while 2-D convolutions operate on
time-frequency representations such as spectrograms and scalograms
\cite{wen2018new,zhang2020deep,zhang2019resnet,hoang2019survey}.
Variants with dilated kernels, multi-scale receptive fields, and residual
connections have progressively improved accuracy under clean-signal
conditions and cross-domain scenarios~\cite{guo2020deep,zhuyun2020cscoh,li2020systematic,jiao2020cnn_review}.

Attention mechanisms and Transformer encoders were subsequently applied
to bearing diagnosis, with self-attention providing a principled mechanism
for relating fault-periodic harmonics across long
windows~\cite{shao2021transformer_bearing,chen2021vibration_transformer}.
However, the quadratic attention cost limits practical window length,
and the position encodings of vanilla Transformers lack explicit
knowledge of bearing kinematics.
Hybrid CNN-attention architectures and attention-based transfer learning
methods have partially addressed these
limitations~\cite{li2022hybrid_fault,liu2024dat}.
Comprehensive surveys of deep learning and transfer learning for bearing
diagnosis are provided in~\cite{yang2023deep_tl,tama2022review}.

A parallel line of work incorporates \emph{physical prior knowledge}
at the signal-processing level, using envelope analysis,
EEMD, and STFT to pre-extract fault-frequency features before feeding
them to classifiers~\cite{randall2011rolling,smith2015rolling}.
Our approach instead embeds kinematic knowledge as a soft regulariser
within end-to-end training, preserving the flexibility of raw-waveform
learning while providing physically motivated frequency alignment.

\subsection{State-Space Models and Mamba for Fault Diagnosis}
\label{subsec:rel_ssm}

Structured state-space models (SSMs)~\cite{gu2021efficiently} — including
S4, S5, H3, and Mamba — have emerged as competitive alternatives to Transformers
for long-sequence modelling, combining linear-time complexity with rich
temporal representations.
Mamba~\cite{gu2023mamba} introduced input-dependent selective state updates,
and Mamba-2~\cite{mamba2} further simplified the formulation by relating
SSM layers to multi-head linear attention.
Mamba-3~\cite{mamba3_2026} generalises Mamba-2 by replacing scalar real states
with complex-valued oscillatory states governed by per-step angle activations,
enabling explicit frequency-domain modelling in the state update.

The efficiency of Mamba on long sequences~\cite{gu2023mamba,mamba2} has
been validated across diverse time-series tasks~\cite{mamba_fault_2024},
and the foundational HiPPO memory projection underpinning SSMs has been
shown to capture long-range dependencies robustly~\cite{mamba_fault2_2025}.
Prior bearing SSM work introduced BearMamba-2 (Mamba-2 backbone,
feature-level noise-awareness, submitted) and KG-cSSM (complex-state priors
at the architecture level for variable-speed diagnosis, under review).
Concurrently, Li~\textit{et~al.}~\cite{li2026pgtmt} propose PG-TMT,
which embeds physics guidance via frequency-domain attention alignment with
bearing fault-order bands, achieving sub-1\,MB edge deployment and EVT-based
reliability calibration; their work positions on operational reliability
and deployment efficiency rather than within-condition accuracy
maximisation, and avoids direct comparison against classical 1D-CNN
baselines.
In contrast, the present work introduces physics directly at the SSM
state level: the kinematic regulariser $\mathcal{L}_{\text{kin}}$ acts on
Mamba-3's complex eigenvalues, which encode oscillation frequency through
their phase angle.
These represent two complementary physics-guidance
strategies---\emph{soft post-hoc alignment} versus \emph{hard
architectural embedding}---targeting different deployment regimes
(edge-IoT versus industrial PC with cross-condition robustness).
The present work is the first to: (i) employ the official Mamba-3 backbone
in a bearing diagnostic context; (ii) exploit the instantaneous angle
activations of Mamba-3 states as the hook for a per-sample kinematic
loss function.

\subsection{Physical Inductive Biases in Neural Diagnostics}
\label{subsec:rel_physics}

Physics-informed neural networks (PINNs)~\cite{raissi2019physics,willard2022physics}
impose governing equations as soft constraints in the training loss, typically
acting on the network's outputs.
In fault diagnosis, physical constraints have been applied at various levels:
physics-inspired correlation fusion integrating domain knowledge into feature
alignment~\cite{chen2022physics_diag},
noise-robust layerwise feature learning with structured constraints on the
loss landscape~\cite{spectrum_reg_2023},
and domain-adversarial ensemble learning that penalises domain-inconsistent
intermediate representations~\cite{complexval_fault}.

Most closely related to our work, Jiang et al.~\cite{jiang2025mechanism}
proposed a mechanism-driven domain adaptation model in which bearing
fault physics are embedded as inductive biases to guide feature alignment
across operating conditions, simultaneously improving cross-domain accuracy
and interpretability.
Our $\mathcal{L}_{\text{kin}}$ shares the spirit of making physical
mechanism knowledge a first-class constraint, but differs in two key
respects: (i) it operates at the \emph{activation level} of a Mamba-3
state-space model — aligning intermediate \emph{instantaneous oscillatory
state frequencies} with analytically known fault harmonics — rather than
at the output or domain-alignment level; (ii) it is architecturally
bound to the Mamba-3 backbone (Section~\ref{subsec:lkin}), making
the backbone selection and physical prior design a single integrated
decision rather than two independent components.
This activation-level operation distinguishes $\mathcal{L}_{\text{kin}}$
from output-level PINNs~\cite{raissi2019physics,willard2022physics}
(which guide predictions, not intermediate features) and from weight-level
or loss-level regularisers (which constrain parameters or domain alignment,
not per-state frequency activations).
Unlike prior work that attracted learned features toward a single fault
class~\cite{lu2017hierarchical}, our cover variant generalises to all
fault frequency families and to variable-speed per-sample targets.

\subsection{Multi-sensor Fusion in Rotating Machinery Diagnosis}
\label{subsec:rel_multisensor}

Early sensor fusion for fault diagnosis concatenated multiscale convolutional
features from multiple sensor locations before a
classifier~\cite{fusion_early,multisensor_encoder}.
More recent approaches fuse at the waveform level through sparse autoencoder
and deep belief network architectures~\cite{chen2017multisensor,wang2020vibro},
attention-based channel alignment~\cite{wan2023multisensor,multisensor_attn},
or multi-sensor residual convolutional fusion under noisy
conditions~\cite{ye2024mrcfn}.
Multi-sensor 2-D CNN fusion has demonstrated accuracy improvements when
sufficient multi-position data are available~\cite{wang2021multisensor2d}.

Orthogonal-axis accelerometers (horizontal and vertical) have been shown
to provide complementary fault signatures for different fault
types~\cite{xjtu_paper}: inner-race faults excite strong vibrations in both
directions, while outer-race faults may be more pronounced in a single
directional axis depending on load zone alignment.
This physical asymmetry directly motivates our finding
(Section~\ref{subsubsec:xjtu}) that dual-sensor fusion benefits IR recall
much more than OR recall in leave-one-bearing-out evaluation.

The precision--sample-efficiency trade-off we characterise
— additional sensor channels improve performance when training data is
sufficient but induce negative transfer in severe few-shot regimes
— has been observed in multi-sensor graph-network studies~\cite{fusion_graph}
and adaptive viewable-graph fusion under adverse conditions~\cite{multisensor_tradeoff},
but has not been systematically documented for SSM-based bearing diagnostics.

\subsection*{Summary and Research Gaps}
\label{subsec:gaps}

In summary, the existing literature reveals two open problems addressed
by this work.
First, while SSMs with complex-valued states carry instantaneous oscillatory
frequencies that have natural physical semantics in bearing diagnostics,
no prior work has exploited these \emph{activation-level} frequencies as
a hook for per-sample kinematic regularisation — existing physics-informed
approaches operate at the output, domain-alignment, or weight level.
Second, the interaction between SSM state bandwidth and multi-sensor input
count — and its dependence on available training data — has not been
characterised.
BearMamba-3 addresses both gaps through an integrated design in which the
Mamba-3 backbone selection and the kinematic loss are structurally bound,
and through systematic single- versus dual-sensor ablations across three
public datasets.
